\scanpage{115}
\deleted{En fait, on demande qu'}En appliquant $\tau_\omega\unsure{=|:|\simeq|\text{à}}\mathbb R^*\tau_\omega$, on trouve
\diagram{\(\text{(2)}\quad\)\begin{tikzcd}[ampersand replacement=\&,column sep=small]
 \deleted{\mathbb R^*}\tau_\omega(\mathbb R^*\Gamma(\mathscr X,S_k))\arrow[r] \& \deleted{\unsure{\mathbb R^*}}\tau_\omega\mathbb R^*\Gamma(\mathscr X,C^{**}).\\
 \& \tau_\omega C^{**}(\mathscr X)\arrow[u]
\end{tikzcd}}
On \deleted{demande} \inserted{voudrait} que ces deux flèches soient des isomorphismes \inserted{dans la cat. dérivée}, ce qui signifie aussi que les hom. déduits de (1) en passant à la coh.
\diagram{\(\text{(1$'$)}\quad\)\begin{tikzcd}[ampersand replacement=\&]
 H^{**}(\mathscr X,S_k)\arrow[r] \& H^{**}(\mathscr X,C^{**})\\
 \&\mathscr H_C^{**}(X)\arrow[u]
\end{tikzcd}}
soient à noyau et conoyau négligeables, ce qui \deleted{signifie} \inserted{implique} précisément, pour la flèche horizontale, qu'on a
\diagram{\(\displaystyle (\deleted{\unsure{1}}2')\left\{\vcenter{\hbox{\begin{tikzcd}[ampersand replacement=\&,column sep=small]
 \tau_\omega(\deleted{\unsure{t}}H^*(\mathscr X,k)\otimes_k S_k)\arrow[r,"\sim"]\arrow[dr,dashed,no head] \& \tau_\omega \mathbb H^{**}(\mathscr X,C^{**})\\
 \&\tau_\omega(\mathbb H_C^{**}(\unsure{X|x}))\arrow[u,"\unsure{\wr|2}"']
\end{tikzcd}}}\right.\) }
% archive sheet number: 114
